\scanpage{321}
and we have
\[
(38)\qquad (-\delta\unsure{v|\nu})B=ac+bc-\unsure{v|\nu}bd
\]
which gives $\uline{u}$ again in terms of $a,b,c,d$, \uline{thanks to} the formula $\uline{\overline{\mathrm{IV}}}(48)$ (page $453'$)
\[
(39)\qquad \uline{u}=\begin{pmatrix}\dfrac1D&\inserted{B/D}\deleted{\unsure{\underline u=}}\\0&1\end{pmatrix}
=\begin{pmatrix}
\dfrac{-\unsure{v|\nu}(ad-bc)}{-\unsure{v|\nu}d^2+cd+d^2}&\dfrac{-ac+bc-\unsure{v|\nu}bd}{-\unsure{v|\nu}\unsure{d|a}^2+cd+d^2}\\0&1
\end{pmatrix}.
\]

\uline{Remark}. On the other hand, the datum of $\beta$ \inserted{\uline{alone}} \inserted{(without the ``sign'' $\varepsilon$)} does not allow one to recover $\mu$. But the datum of $(\beta,\unsure{\varepsilon})$ does, for from it one deduces the knowledge of
\[
(\unsure{40})\qquad \eta=\varepsilon\beta\sigma(\beta)^{-1}\overset{\mathrm{dfn}}=\begin{pmatrix}R&S\\T&U\end{pmatrix},
\]
which will be related to $\zeta=\alpha\rho(\alpha)^{-1}$ by
\[
\zeta=\eta\varepsilon_1(\mu-1)
\]
and by the theory \unsure{of the|of} (II) we will have
\[
(41)\qquad \mu=1/U
\]
(but in terms of $\beta\sigma(\beta)^{-1}\overset{\mathrm{dfn}}=\begin{pmatrix}R_0&S_0\\T_0&U_0\end{pmatrix}$,
we will have \deleted{\illegible{1}} \unsure{)}
\[
(42)\qquad \mu=\frac{\varepsilon}{U_0}\qquad)
\]
i.e. it is \inserted{a priori} not known from $U_0$ alone, but only via $U_0$ and $\varepsilon$. \deleted{\unsure{R}}

To make this point precise, we are going to look at \inserted{when and} how, for fixed $(\beta,\varepsilon)$, one can complete with $\alpha,\mu$ so as to have a
\mspage{320}

\scanpage{322}
\mspage{461}
section $(\alpha,\beta,\mu,\varepsilon)$ of $\mathcal G^{*}_{\rho,\sigma}$. A necessary condition is that $U_0$ be invertible (it depends only on $\beta$, not on $\varepsilon$) --- and then one can choose $\varepsilon$ arbitrarily in $\Gamma\mu_2$, \deleted{and}\inserted{$\varepsilon\in\Gamma\mu_2$} $\eta=\varepsilon\beta\sigma(\beta)^{-1}$ will then be a section of the \inserted{\unsure{subscheme}} $A_\sigma$ \deleted{\illegible{1}} of $G$ of solutions of $\det\eta=1$, $\operatorname{Tr}\eta\sigma=\operatorname{Tr}\sigma$ (cf. lemma page 455), hence by the theory of (II), for \deleted{\unsure{completing taking}} $\mu=1/U=\varepsilon/U_0$, \deleted{\illegible{2}} and defining
\[
\zeta=\eta\varepsilon_1(\mu-1),
\]
$\zeta$ will indeed be expressible \inserted{locally} in the form $\alpha\rho(\alpha)^{-1}$, the indeterminacy being moreover in $T_\rho$ via $\alpha\longmapsto\alpha'=\alpha\uline b$, $\uline b\in\Gamma T_\rho$.

This does show that $\mu$ can \unsure{thereof|not} take, for a fixed $\beta$, all the values $\varepsilon/U_0$, where $\beta\sigma(\beta)^{-1}=\begin{pmatrix}R_0&\unsure{S_0|C_0}\\T_0&U_0\end{pmatrix}$ \ldots
\mspage{321}

\scanpage{323}
\mspage{462}
($\uline{\overline{\mathrm{VI}}}$) Characterization \unsure{of the|of the} modular triple $(\varepsilon_1,\rho,\sigma)$.

$\unsure{1^\circ})$ Let $G,SG,PG$ be associated forms of $\mathrm{Gl}(2),\mathrm{Sl}(2),\mathrm{GP}(1)$ over a base $S$ --- thus corresponding: a torsor under $\mathrm{GP}(1)_S$, or again: a bundle of projective lines \inserted{$P$} over $S$, or an Azumaya Algebra $M$ of rank $\unsure{1|4}$, and let us consider two sections \inserted{\illegible{1}} $\widetilde\rho,\widetilde\sigma$ of $PG$, such that
\[
(1)\qquad \unsure{a)}\quad\widetilde\varepsilon_1=\widetilde\rho\widetilde\sigma^{-1}
\]
be \deleted{\illegible{1}} a \uline{regular} unipotent section, and that the section
\[
(2)\qquad \widetilde\varepsilon_0\overset{\mathrm{dfn}}=\widetilde\rho^{-1}(\widetilde\varepsilon_1)=\widetilde\sigma^{-1}(\widetilde\varepsilon_1)
\]
which is deduced from it be \deleted{\unsure{everywhere}} \inserted{in each fiber} inequivalent to $\widetilde\varepsilon_0$ (i.e. that the sections $x_0,x_1$ of $P$ which correspond to $\varepsilon_0,\varepsilon_1$ be \deleted{\unsure{everywhere}} distinct in each fiber).

\margin{\uline{N.B.} If $\varepsilon_0,\varepsilon_1$ are the canonical liftings of $\widetilde\varepsilon_0,\widetilde\varepsilon_1$ into $\unsure{G|P}$, \unsure{regular}, \inserted{with $\unsure{\widetilde\sigma(\varepsilon_0)=\widetilde\rho(\varepsilon_0)=\varepsilon_1|\widetilde\sigma(\widetilde\varepsilon_0)=\widetilde\rho(\widetilde\varepsilon_0)=\widetilde\varepsilon_1}$} one can lift $\widetilde\rho,\widetilde\sigma$ in a \uline{unique} way to $\rho,\sigma$ so as to satisfy
$-\sigma^{-1}\rho=\varepsilon_0\quad(\Longleftrightarrow-\rho\sigma^{-1}=\varepsilon_1)$,
$(\rho+\sigma)^2=(\rho+\sigma)$.}

Then there exists a \uline{unique}$^*$ pinning
\[
(3)\qquad\deleted{\unsure{G_0=}}\quad\mathrm{Gl}(2)_S\xrightarrow[\sim]{\varphi}G
\]
and sections $\unsure{v|\nu},\unsure{w|\omega},\unsure{v|\nu}',\unsure{w|\omega}'$ of $\mathcal O_S$, satisfying
\[
(4)\qquad \unsure{v|\nu}+\unsure{v|\nu}'=0,\qquad \unsure{w|\omega}+\unsure{w|\omega}'=1,\qquad \unsure{v|\nu}\in\Gamma\mathcal O_S^*
\]
such that $\widetilde\rho,\widetilde\sigma$ are the images, under the composite $\mathrm{Gl}(2)_S\xrightarrow[\sim]{\varphi}G\longrightarrow PG$, of the sections
\[
(5)\qquad \rho=\begin{pmatrix}0&1\\\unsure{v|\nu}&\unsure{w|\omega}\end{pmatrix},\qquad
\sigma=\begin{pmatrix}0&-1\\\unsure{v|\nu}'&\unsure{w|\omega}'\end{pmatrix}
\]
\margin{\uline{N.B.}$^*$ Uniqueness \inserted{of $\unsure{v|\nu},\unsure{w|\omega}$ (i.e. of $\rho_0,\sigma_0$)} \unsure{is} not true without some restrictions on $S$, it suffices that every section $\tau$ of $\mathcal O_S$ such that $2\tau=0,\ \tau^2=0$ be zero.}
\mspage{322}

\scanpage{324}
of $\mathrm{Gl}(2)_S$ --- and conversely, for every pinning $\varphi$ and \unsure{every system} \deleted{\illegible{4}}\inserted{$\unsure{v|\nu},\unsure{w|\omega},\unsure{v|\nu}',\unsure{w|\omega}'$} satisfying (4) --- i.e.
\[
(4\mathrm{bis})\qquad \unsure{v|\nu}'=-\unsure{v|\nu},\quad \unsure{w|\omega}'=1-\unsure{w|\omega},\quad\begin{gathered}\unsure{v|\nu}\in\Gamma(S,\mathcal O_S^*)\\\unsure{w|\omega}\in\Gamma(S,\mathcal O_S)\end{gathered},
\]
the images \inserted{$\widetilde\rho,\widetilde\sigma$} under the composite $\mathrm{Gl}(2)_S\longrightarrow PG\simeq G/\mathbf G_m$ \deleted{satisfy}\inserted{of the} sections $\rho,\sigma$ of $\mathrm{Gl}(2)_S$ given by (5) satisfy the conditions above ($\widetilde\rho\widetilde\sigma^{-1}\overset{\mathrm{dfn}}=\widetilde\varepsilon_1$ is \uline{regular} unipotent, and everywhere inequivalent to $\widetilde\sigma^{-1}(\widetilde\varepsilon_1)=\widetilde\rho^{-1}(\widetilde\varepsilon_1)\overset{\mathrm{dfn}}=\widetilde\varepsilon_0$).

Setting then
\[
(6)\qquad \varepsilon_{\unsure{1|0}}=\begin{pmatrix}1&0\\1&1\end{pmatrix},\qquad \varepsilon_0=\begin{pmatrix}1&1/\unsure{v|\nu}\\0&1\end{pmatrix}
\]
we will have
\[
(7)\qquad \rho\sigma^{-1}=-\varepsilon_1,\qquad\sigma^{-1}\rho=-\varepsilon_0
\]
\[
(7\mathrm{bis})\qquad \varepsilon_0=\rho^{-1}(\varepsilon_1)=\sigma^{-1}(\varepsilon_1),\qquad\varepsilon_1=\rho(\varepsilon_0)=\sigma(\varepsilon_0)
\]
\uline{whence} the images under $\mathrm{Gl}(2)_S\longrightarrow PG$ of \deleted{$\rho$ and $\sigma$} $\varepsilon_1,\varepsilon_0$ are $\widetilde\varepsilon_1,\widetilde\varepsilon_0$. We have the formulas
\[
(8)\qquad \operatorname{Tr}\sigma=\unsure{w|\omega}',\quad\operatorname{Tr}\rho=\unsure{w|\omega},\quad\operatorname{Tr}(\sigma+\rho)=1
\]
\[
(9)\qquad \sigma+\rho=\begin{pmatrix}0&0\\0&1\end{pmatrix}\overset{\mathrm{dfn}}=N\deleted{\unsure{0}}
\]
\[
(10)\qquad \operatorname{Tr}\zeta N=D\quad\text{if}\quad\zeta=\begin{pmatrix}A&B\\C&\uline D\end{pmatrix},
\]
\[
(11)\qquad \det\sigma=\det\rho=-\unsure{v|\nu}=\unsure{-v'|v'}.
\]
\mspage{323}

\scanpage{325}
\mspage{463}
$2^\circ)$ The following conditions are equivalent
\diagram{\[
(12)\quad\left\{
\begin{tikzcd}[ampersand replacement=\&,column sep=large,row sep=large]
\boxed{\operatorname{Tr}\sigma=0}\arrow[r,Leftrightarrow]\arrow[d,Leftrightarrow]\&\operatorname{Tr}\rho=1\arrow[r,Leftrightarrow]\arrow[d,Leftrightarrow]\&\boxed{\widetilde\sigma^2=1}\arrow[d,Leftrightarrow]\\
\deleted{\mathrm{Tr}}\boxed{\unsure{w|\omega}'=0}\&\boxed{\unsure{w|\omega}=1}\&\sigma^2\in\Gamma\mathbf G_m\arrow[d,Leftrightarrow]\\
\&\&\boxed{\sigma^2=\unsure{v|\nu}\,\mathrm{id}.}
\end{tikzcd}\right.
\]}
They are also equivalent to this: denoting by $N'_\sigma$ the subgroup scheme of $\mathrm{Gl}(2)$ defined by the condition
\[
(13)\qquad \tau(\widetilde\sigma)=\widetilde\sigma\quad\text{i.e.}\quad\tau(\sigma)=\varepsilon\sigma,\quad\varepsilon\in\Gamma\mathbf G_m
\]
(we will necessarily have $\varepsilon^2=1$ i.e. $\varepsilon\in\Gamma\mu_2$)
the morphism
\[
(14)\qquad\begin{aligned}N'_\sigma&\longrightarrow\mu_2\\\tau&\longmapsto\varepsilon\end{aligned}
\]
is an \uline{epimorphism} --- so that we will have an exact sequence
\[
(15)\qquad1\longrightarrow T_\sigma\longrightarrow N'_\sigma\longrightarrow\mu_2\longrightarrow1,
\]
where \deleted{\unsure{Let us set}}
\[
(16)\qquad T_\sigma=\operatorname{Centr}_{\mathrm{Gl}(2)_S}(\sigma)=\{t\sigma+u\mid\det(t\sigma+u)\ \text{invertible}\}
\]
\[
(17)\qquad\text{We have}\quad N'_\sigma\subset N_\sigma\overset{\mathrm{dfn}}=\operatorname{Norm}_{\mathrm{Gl}(2)}(T_\sigma).
\]
If $2$ is a regular section of $\mathcal O_S$, then every section of $N_\sigma$ \inserted{over $S$} is a section of $N'_\sigma$ --- thus \unsure{if|in} \inserted{the} residue characteristics \unsure{are|is} $\ne2$ i.e. $2$ invertible, we will have
\[
(18)\qquad N'_\sigma=N_\sigma\qquad(\text{residue characteristics }\ne2).
\]
The condition under consideration means that $\sigma,\rho$ take the form
\mspage{324}

\scanpage{326}
\[
(19)\qquad\sigma=\begin{pmatrix}0&-1\\\underset{\substack{\Vert\\-\unsure{v|\nu}}}{\unsure{v|\nu}'}&0\end{pmatrix},\qquad\rho=\begin{pmatrix}0&1\\\unsure{v|\nu}&1\end{pmatrix}
\]
with $\unsure{v|\nu}+\unsure{v|\nu}'=0$.

These results of $1^\circ),2^\circ)$ are only a recapitulation of part of the results of the preceding sections $\uline{\overline{\mathrm{IV}}},\uline{\overline{\mathrm{V}}}$. \deleted{\unsure{Here}}, we have just studied conditions which amount to the determination \inserted{\unsure{the value of|the values of}} $\unsure{w|\omega}$ (or again \inserted{of} $\unsure{w|\omega}'$) in a \deleted{quite} explicit way. We are going to \unsure{turn to} conditions which \unsure{are to|will} serve to fix $\unsure{v|\nu}$ (or again $\unsure{v|\nu}'$) in an explicit way.

$3^\circ)$ Equivalent conditions
\[
(20)\quad\left\{\begin{aligned}
\det\sigma=1&\Longleftrightarrow\det\rho=1\Longleftrightarrow \unsure{v|\nu}=-1\Longleftrightarrow \unsure{v|\nu}'=1\\
&\Longleftrightarrow\varepsilon_0=\begin{pmatrix}1&-1\\0&1\end{pmatrix}.
\end{aligned}\right.
\]
These equivalences follow at once from (6) and (11)\unsure{.|!} The conditions also mean that $\rho$ and $\sigma$ have the form
\[
(21)\qquad \deleted{\rho}\inserted{\sigma}=\begin{pmatrix}0&-1\\1&\unsure{w|\omega}'\end{pmatrix},\qquad\rho=\begin{pmatrix}0&1\\-1&\unsure{w|\omega}\end{pmatrix}
\]
\[
(\unsure{w|\omega}+\unsure{w|\omega}'=1)
\]
$4^\circ)$ Conjunction of $2^\circ)$ and $3^\circ)$.

Equivalent conditions:
\mspage{325}

\scanpage{327}
\mspage{464}

\[
(22)\quad
\left\{
\begin{array}{ccccc}
\begin{gathered}
\text{Conjunction of the}\\
\text{conditions of }2^\circ\\
\text{and of the conditions}\\
\text{of }3^\circ
\end{gathered}
&\Longleftrightarrow&
\begin{gathered}\boxed{\sigma^2=-1}\\
\text{i.e. }\sigma^2+1=0\end{gathered}
&\Longleftrightarrow&
\boxed{\widetilde\sigma^2=\widetilde\rho^3=1}
\\
\Updownarrow&&&&\\
\unsure{v|\nu}=-1,\unsure{w|\omega}=1&\Longleftrightarrow&
\rho=\begin{pmatrix}0&1\\-1&1\end{pmatrix}
&\Longleftrightarrow&\operatorname{Tr}\rho=\det\rho=1
\\
\Updownarrow&&&&\\
\unsure{v|\nu}'=1,\unsure{w|\omega}'=0&\Longleftrightarrow&
\sigma=\begin{pmatrix}0&-1\\1&0\end{pmatrix}
&\Longleftrightarrow&(\operatorname{Tr}\sigma=0,\det\sigma=1)
\end{array}\right.
\]
\margin{\diagram{\begin{tikzcd}[
ampersand replacement=\&,
column sep=large,
row sep=small,
execute at end picture={\draw[-] (quadratic.north east) -- (quadratic.north west) -- (quadratic.south west) -- (quadratic.south east);}
]
\boxed{\widetilde\sigma^2=\widetilde\rho^3=1}
\arrow[r,Leftrightarrow]
\arrow[dr,Leftrightarrow]
\& \deleted{\boxed{\rho^3=-1}}
\\
\& |[alias=quadratic,inner sep=3pt]| {\rho^2-\rho+1=\unsure{\illegible{1}}}
\end{tikzcd}}
cf. corollary~2 below}

\uline{Proof.} The fact that the conjunction of the conditions of $2^\circ)$ and $3^\circ)$ is equivalent to \unsure{each} of the conditions of the $2^{\mathrm e}$ and $3^{\mathrm e}$ rows, is trivial. The equivalence with $\sigma^2=-1$ comes from the fact that \deleted{\illegible{1}} the condition $2^\circ$ means that \inserted{$\sigma^2$ \unsure{be|is} scalar, i.e.} $\sigma^2=\unsure{v|\nu}\,\mathrm{id}$, \deleted{\illegible{2}} whereas $3^\circ$ means $\unsure{v|\nu}=-1$, so $\sigma^2=-1$ is indeed equivalent to the conjunction.

\margin{for another proof of the equivalence of $2^\circ)+3^\circ)$ with $\widetilde\sigma^2=\widetilde\rho^3=1$, cf. lemma below}

Let us note that $2^\circ)$ is also written $\widetilde\sigma^2=1$, and is equivalent to: $\unsure{w|\omega}'=0$, i.e. $\unsure{w|\omega}=1$, and we then find, writing \deleted{\unsure{$\rho=$}} 
\[
\rho=\begin{pmatrix}0&1\\\unsure{v|\nu}&1\end{pmatrix},
\]
\[
(23)\qquad
\rho^3=\begin{pmatrix}\unsure{v|\nu}&\unsure{v|\nu}+1\\\unsure{v|\nu}(\unsure{v|\nu}+1)&2\unsure{v|\nu}+1\end{pmatrix}
\]
thus $\rho^3$ can be a scalar i.e. $\widetilde\rho^3=1$ only if the two diagonal terms are equal, now their difference is $(2\unsure{v|\nu}+1)-\unsure{v|\nu}=\unsure{v|\nu}+1$,

\mspage{326}